\documentclass[12pt,a4paper]{article}

% ============================================================
% Drake Journal Reading LaTeX Template
% 作者: 謝慕揚醫師 MD, PhD, FESC
% ============================================================

% Chinese font setup
\usepackage{xeCJK}
\usepackage{fontspec}
\IfFontExistsTF{Noto Serif CJK TC}{
  \setCJKmainfont{Noto Serif CJK TC}
  \setCJKsansfont{Noto Serif CJK TC}
  \setCJKmonofont{Noto Serif CJK TC}
}{\setCJKmainfont{SimSun}
  \setCJKsansfont{SimSun}
  \setCJKmonofont{SimSun}
}

% Essential packages
\usepackage[margin=2cm,headheight=25pt,top=2.5cm]{geometry}
\usepackage{graphicx}
\usepackage{booktabs}
\usepackage{longtable}
\usepackage{array}
\usepackage{xcolor}
\usepackage{hyperref}
\usepackage{tcolorbox}
\usepackage{titlesec}
\usepackage{enumitem}
\usepackage{amsmath}
\usepackage{multirow}
\usepackage{fancyhdr}
\usepackage{datetime2}

% Color definitions
\definecolor{emergency_red}{RGB}{186,24,27}
\definecolor{emergency_blue}{RGB}{0,114,188}
\definecolor{emergency_gray}{RGB}{120,120,120}
\definecolor{highlight_yellow}{RGB}{255,250,205}

% Hyperlink setup
\hypersetup{
    colorlinks=true,
    linkcolor=emergency_blue,
    citecolor=emergency_blue,
    urlcolor=emergency_blue,
    pdfborder={0 0 0}
}

% Title formats
\titleformat{\section}{\Large\bfseries\color{emergency_red}}{\thesection}{1em}{}
\titleformat{\subsection}{\large\bfseries\color{emergency_blue}}{\thesubsection}{1em}{}
\titleformat{\subsubsection}{\normalsize\bfseries\color{emergency_gray}}{\thesubsubsection}{1em}{}

% Header/Footer setup
\pagestyle{fancy}
\fancyhf{}
\fancyhead[L]{\textcolor{emergency_red}{NEJM AI 教學講義}}
\fancyhead[R]{\textcolor{emergency_blue}{Medical LLM Agents}}
\fancyfoot[C]{\textcolor{emergency_gray}{\thepage}}
\renewcommand{\headrulewidth}{1pt}
\renewcommand{\headrule}{\hbox to\headwidth{\color{emergency_red}\leaders\hrule height \headrulewidth\hfill}}

% Date format - use existing iso style
\DTMsetdatestyle{iso}
\newcommand{\mydate}{\DTMtoday}

\begin{document}

% ============================================================
% Title Page
% ============================================================

\begin{center}
{\LARGE\bfseries\textcolor{emergency_red}{NEJM AI}}\\[0.5cm]
{\Large\textbf{MedAgentBench: A Virtual EHR Environment to Benchmark Medical LLM Agents}}\\[0.3cm]
{\large MedAgentBench：評估醫療大型語言模型代理的虛擬電子病歷環境}\\[0.5cm]
{\normalsize NEJM AI 2025;2(9). DOI: 10.1056/AIdbp2500144}\\[0.3cm]
{\normalsize 整理：謝慕揚 MD, PhD, FESC \quad 日期：\mydate}
\end{center}

\vspace{0.5cm}

% Key findings box
\begin{tcolorbox}[colback=yellow!10!white,colframe=orange!75!black,title=\textbf{核心發現摘要}]
\textbf{重大發現：}本研究建立首個評估 LLM 在醫療記錄環境中代理能力的標準化基準。目前最佳模型 Claude 3.5 Sonnet v2 成功率為 69.67\%，顯示 AI 代理在醫療應用中具有潛力，但仍需改進。

\textbf{臨床意義：}醫師僅 27\% 時間用於直接臨床照護，其餘為文書與行政工作。AI 代理可望減輕行政負擔，協助醫師回歸臨床照護。
\end{tcolorbox}

\tableofcontents
\newpage

% ============================================================
\section{研究概述}
% ============================================================

本研究來自 Stanford University，建立了 MedAgentBench 評估套件，用於測試大型語言模型 (Large Language Models, LLMs) 在醫療記錄環境中的代理能力。這是首個需要與醫療記錄環境進行自主互動的標準化基準。

\subsection{從 Chatbot 到 Agent 的演進}

\begin{itemize}
    \item \textbf{傳統 LLM (Chatbot)：}依賴使用者提示，提供獨立輸出
    \item \textbf{AI Agent：}主動解讀高層指令、規劃行動、與外部系統互動、迭代優化回應
    \item \textbf{關鍵轉變：}從「AI 作為工具」轉向「AI 作為隊友」
\end{itemize}

\begin{tcolorbox}[colback=emergency_blue!5!white,colframe=emergency_blue,title=\textbf{Chatbot vs Agent 範例}]
\textbf{臨床情境：}社區型肺炎 (Community-Acquired Pneumonia, CAP) 治療

\textbf{Chatbot 回應：}回答「CAP 的住院治療方案是什麼？」

\textbf{Agent 執行：}
\begin{enumerate}
    \item 計算個人化病人風險評分
    \item 評估 Pseudomonas 感染風險因子
    \item 分析藥物交互作用與過敏史
    \item 整合先前培養數據與本地抗生素敏感性資料
    \item 排隊抗生素與支持性照護醫囑供醫師審核簽署
    \item 自主排程追蹤門診
\end{enumerate}
\end{tcolorbox}

% ============================================================
\section{研究資訊}
% ============================================================

\subsection{基本資料}

\textbf{標題：}MedAgentBench: A Virtual EHR Environment to Benchmark Medical LLM Agents\\
\textbf{作者：}Yixing Jiang, Kameron C. Black DO MPH, Gloria Geng, Danny Park, James Zou PhD, Andrew Y. Ng PhD, Jonathan H. Chen MD PhD\\
\textbf{機構：}Stanford University\\
\textbf{期刊：}NEJM AI 2025;2(9)\\
\textbf{DOI：}\href{https://doi.org/10.1056/AIdbp2500144}{10.1056/AIdbp2500144}\\
\textbf{資源：}\href{https://github.com/stanfordmlgroup/MedAgentBench}{https://github.com/stanfordmlgroup/MedAgentBench}

\subsection{研究背景}

\begin{itemize}
    \item 醫師僅約 \textbf{27\%} 時間用於直接臨床照護，其餘用於文書與行政工作
    \item 傳統 QA-based AI 基準 (如 MedQA、MedMCQA) 已達飽和，模型表現已達超人水準
    \item 缺乏評估 LLM 在醫療情境中代理能力的標準化基準
    \item 需要基於 FHIR (Fast Healthcare Interoperability Resources) API 的基準測試
\end{itemize}

% ============================================================
\section{MedAgentBench 框架}
% ============================================================

\subsection{系統架構}

\begin{enumerate}
    \item \textbf{臨床醫師指定任務：}例如「檢查病人最近的鉀離子數值，若過低則開立補充」
    \item \textbf{Agent 協調器：}解讀任務並規劃 function calls
    \item \textbf{執行任務：}透過 FHIR server 發送請求修改醫療記錄資料庫
    \item \textbf{回饋給臨床醫師：}摘要已執行的任務
\end{enumerate}

\subsection{任務類別}

\begin{longtable}{p{3.5cm}p{5cm}p{6cm}}
\caption{MedAgentBench 任務類別總覽} \\
\toprule
\textbf{任務類別} & \textbf{範例指令} & \textbf{EHR 系統情境} \\
\midrule
\endfirsthead

\toprule
\textbf{任務類別} & \textbf{範例指令} & \textbf{EHR 系統情境} \\
\midrule
\endhead

\bottomrule
\endfoot

病人資訊檢索\newline (Patient Information Retrieval) &
查詢姓名與生日對應的 MRN &
N/A \\
\midrule

檢驗結果檢索\newline (Laboratory Result Retrieval) &
查詢過去 24 小時最近的 Magnesium level &
提供時間、檢驗代碼、單位換算規則 \\
\midrule

病人數據彙總\newline (Patient Data Aggregation) &
計算過去 24 小時的平均血糖值 &
提供時間與 GLU 代碼 \\
\midrule

記錄病人數據\newline (Recording Patient Data) &
記錄血壓測量值 118/77 mmHg &
提供時間與 flowsheet ID \\
\midrule

檢查開立\newline (Test Ordering) &
查詢 HbA1C，若超過一年則重新開立 &
提供 LOINC code: 4548-4 \\
\midrule

轉介開立\newline (Referral Ordering) &
開立骨科手術轉介 &
提供 SNOMED code \\
\midrule

藥物開立\newline (Medication Ordering) &
檢查 K+ level，若過低則依劑量指示開立補充 &
提供 NDC code、劑量計算規則、LOINC code \\
\bottomrule
\end{longtable}

\subsection{病人資料}

\begin{longtable}{p{6cm}c}
\toprule
\textbf{特徵} & \textbf{數值} \\
\midrule
\endfirsthead

\toprule
\textbf{特徵} & \textbf{數值} \\
\midrule
\endhead

\bottomrule
\endfoot

獨立個體數 & 100 \\
年齡 (years, mean±SD) & 58.15±19.82 \\
女性比例 & 47\% \\
總記錄數 & 785,207 \\
Observation records & 563,426 \\
Procedure records & 124,969 \\
Condition records & 74,821 \\
MedicationRequest records & 21,991 \\
\bottomrule
\end{longtable}

% ============================================================
\section{主要結果}
% ============================================================

\subsection{模型表現比較}

\begin{longtable}{p{4cm}ccccc}
\caption{各模型在 MedAgentBench 的表現} \\
\toprule
\textbf{模型} & \textbf{大小} & \textbf{形式} & \textbf{整體 SR} & \textbf{Query SR} & \textbf{Action SR} \\
\midrule
\endfirsthead

\toprule
\textbf{模型} & \textbf{大小} & \textbf{形式} & \textbf{整體 SR} & \textbf{Query SR} & \textbf{Action SR} \\
\midrule
\endhead

\bottomrule
\endfoot

\textbf{Claude 3.5 Sonnet v2} & N/A & API & \textbf{69.67\%} & \textbf{85.33\%} & 54.00\% \\
GPT-4o & N/A & API & 64.00\% & 72.00\% & 56.00\% \\
DeepSeek-V3 & 685B & Open & 62.67\% & 70.67\% & 54.67\% \\
Gemini-1.5 Pro & N/A & API & 62.00\% & 52.67\% & \textbf{71.33\%} \\
GPT-4o-mini & N/A & API & 56.33\% & 59.33\% & 53.33\% \\
o3-mini & N/A & API & 51.67\% & 54.67\% & 48.67\% \\
Qwen2.5 & 72B & Open & 51.33\% & 38.67\% & 64.00\% \\
Llama 3.3 & 70B & Open & 46.33\% & 50.00\% & 42.67\% \\
Gemini 2.0 Flash & N/A & API & 38.33\% & 34.00\% & 42.67\% \\
Gemma2 & 27B & Open & 19.33\% & 38.67\% & 0.00\% \\
Gemini 2.0 Pro & N/A & API & 18.00\% & 25.33\% & 10.67\% \\
Mistral v0.3 & 7B & Open & 4.00\% & 8.00\% & 0.00\% \\
\bottomrule
\end{longtable}

\textbf{註：}SR = Success Rate (成功率)；Query SR = 查詢型任務成功率；Action SR = 行動型任務成功率

\subsection{關鍵發現}

\begin{enumerate}
    \item \textcolor{blue}{\textbf{最佳表現：}}Claude 3.5 Sonnet v2 整體成功率 69.67\%，Query 任務達 85.33\%
    \item \textcolor{blue}{\textbf{Query vs Action：}}大多數模型在查詢型任務表現優於行動型任務
    \item \textcolor{blue}{\textbf{開源 vs 閉源：}}閉源模型仍優於開源模型，開源社群仍有改進空間
    \item \textcolor{blue}{\textbf{臨床應用潛力：}}可優先探索僅需資訊檢索的使用案例
\end{enumerate}

\subsection{常見錯誤模式}

\begin{itemize}
    \item \textbf{未精確遵循指令：}如 Gemini 2.0 Flash 在 54\% 案例中輸出無效動作
    \item \textbf{答案格式錯誤：}模型輸出完整句子，而非預期的純數值
    \item \textbf{範例：}模型輸出 "[``value'': 5.4]"，但預期答案為 "[5.4]"
\end{itemize}

% ============================================================
\section{Agentic AI 在醫療的潛在應用}
% ============================================================

\subsection{可自動化的臨床工作流程}

\begin{longtable}{p{4cm}p{10cm}}
\toprule
\textbf{應用領域} & \textbf{說明} \\
\midrule
\endfirsthead

\toprule
\textbf{應用領域} & \textbf{說明} \\
\midrule
\endhead

\bottomrule
\endfoot

術前風險評估\newline (Preoperative Risk Assessment) &
自動整合病人資料評估手術適應症 \\
\midrule

法規遵循監測\newline (Regulatory Compliance) &
監測醫院安全措施的法規遵循情況 \\
\midrule

臨床分診\newline (Clinical Triage) &
初步評估病人緊急程度與分類 \\
\midrule

EHR 系統設定\newline (EHR Configuration) &
協助電子病歷系統的維護與設定 \\
\midrule

保險事前授權\newline (Prior Authorization) &
自動化保險事前授權申請流程 \\
\midrule

病人溝通\newline (Patient Communication) &
處理病人訊息與回覆 \\
\midrule

文件撰寫\newline (Documentation) &
協助臨床文件撰寫與紀錄 \\
\bottomrule
\end{longtable}

% ============================================================
\section{臨床醫師可現行採用的 AI 工具與協定}
% ============================================================

\begin{tcolorbox}[colback=emergency_blue!5!white,colframe=emergency_blue,title=\textbf{實用工具與建議}]

\subsection*{一、現行可用的 AI 輔助工具}

\begin{enumerate}
    \item \textbf{Ambient Clinical Intelligence (環境臨床智慧)}
    \begin{itemize}
        \item Nuance DAX Copilot (Microsoft/Nuance)
        \item Abridge
        \item Suki AI
        \item 功能：自動將醫病對話轉錄為臨床文件
    \end{itemize}

    \item \textbf{Clinical Documentation AI}
    \begin{itemize}
        \item Epic 整合的 AI 功能 (InBasket message drafting)
        \item 各 EHR 廠商的 LLM 整合功能
        \item 功能：草擬病人訊息回覆、摘要病歷
    \end{itemize}

    \item \textbf{通用 LLM 工具 (需注意 PHI 保護)}
    \begin{itemize}
        \item ChatGPT、Claude、Gemini (不可輸入 PHI)
        \item 功能：文獻搜尋、教育、一般寫作輔助
    \end{itemize}
\end{enumerate}

\subsection*{二、FHIR 標準與互操作性}

\begin{itemize}
    \item \textbf{FHIR (Fast Healthcare Interoperability Resources)：}美國 21st Century Cures Act 規定的醫療資料存取標準
    \item 多數商用 EHR 廠商已支援 FHIR API
    \item 可作為 AI Agent 與 EHR 系統整合的標準介面
\end{itemize}

\subsection*{三、採用建議}

\begin{enumerate}
    \item \textbf{優先探索查詢型任務：}目前 LLM 在資訊檢索任務表現較佳 (Query SR 85.33\% vs Action SR 54\%)
    \item \textbf{保持人工監督：}目前最佳模型成功率僅 69.67\%，仍需醫師審核
    \item \textbf{關注 PHI 保護：}使用任何 AI 工具時需確保病人資料保護
    \item \textbf{選擇經驗證的工具：}優先選用經 FDA/CE 認證或與 EHR 廠商整合的工具
\end{enumerate}

\end{tcolorbox}

% ============================================================
\section{研究限制}
% ============================================================

\begin{enumerate}
    \item \textbf{無法捕捉完整真實情境：}真實醫療情境需要團隊間的協調與溝通
    \item \textbf{缺少臨床筆記：}病人資料中未包含 clinical notes 等結構化數據
    \item \textbf{需要醫院特定設定：}需要 hospital-specific EHR system context
    \item \textbf{病人代表性：}所有病人資料來自 Stanford Hospital，可能存在偏差
    \item \textbf{任務範圍：}專注於醫療記錄情境，未涵蓋外科或護理領域
\end{enumerate}

% ============================================================
\section{結論}
% ============================================================

\begin{enumerate}
    \item \textcolor{red}{\textbf{結論一：}}MedAgentBench 是首個評估 LLM 在醫療記錄環境中代理能力的標準化基準，使用 FHIR 標準可輕鬆遷移至實際 EHR 系統

    \item \textcolor{red}{\textbf{結論二：}}目前最先進的 LLM 展現初步能力 (最佳 69.67\%)，但尚無法可靠處理所有臨床相關任務的完整複雜性

    \item \textcolor{red}{\textbf{結論三：}}LLM 在查詢型任務表現優於行動型任務，建議優先探索資訊檢索類的臨床應用

    \item \textcolor{red}{\textbf{結論四：}}Agent-based task frameworks 是推進 AI 系統有效整合至臨床工作流程的必要下一步
\end{enumerate}

\vspace{0.5cm}

\begin{tcolorbox}[colback=highlight_yellow,colframe=emergency_red,title=\textbf{Take-Home Message}]
\begin{enumerate}
    \item AI Agent 從「工具」演進為「隊友」，可主動規劃與執行臨床任務
    \item 目前技術成熟度：查詢任務較可靠，行動任務仍需改進
    \item 臨床應用建議：優先採用資訊檢索類工具，保持人工監督
    \item FHIR 標準是 AI 與 EHR 整合的關鍵基礎設施
    \item 持續關注 MedAgentBench 進展可了解 AI 醫療應用的最新發展
\end{enumerate}
\end{tcolorbox}

\vspace{0.5cm}
\textbf{關鍵詞：}Large Language Model (LLM)、AI Agent、Electronic Health Record (EHR)、FHIR、Clinical Documentation、MedAgentBench

% ============================================================
\section{參考文獻}
% ============================================================

\begin{enumerate}
    \item Jiang Y, Black KC, Geng G, et al. MedAgentBench: A Virtual EHR Environment to Benchmark Medical LLM Agents. \href{https://doi.org/10.1056/AIdbp2500144}{\textit{NEJM AI}. 2025;2(9).}

    \item Zou J, Topol EJ. The rise of agentic AI teammates in medicine. \href{https://doi.org/10.1016/S0140-6736(25)00202-8}{\textit{Lancet}. 2025;405:457.}

    \item Sahni NR, Carrus B. Artificial intelligence in US health care delivery. \href{https://doi.org/10.1056/NEJMra2204673}{\textit{N Engl J Med}. 2023;389:348-358.}

    \item Sinsky C, Colligan L, Li L, et al. Allocation of physician time in ambulatory practice: a time and motion study in 4 specialties. \href{https://doi.org/10.7326/M16-0961}{\textit{Ann Intern Med}. 2016;165:753-760.}

    \item Tripathi S, Sukumaran R, Cook TS. Efficient healthcare with large language models: optimizing clinical workflow and enhancing patient care. \href{https://doi.org/10.1093/jamia/ocad258}{\textit{J Am Med Inform Assoc}. 2024;31:1436-1440.}
\end{enumerate}

% ============================================================
\section{縮寫對照表}
% ============================================================

\begin{longtable}{p{4cm}p{10cm}}
\toprule
\textbf{縮寫} & \textbf{全名} \\
\midrule
\endfirsthead

\toprule
\textbf{縮寫} & \textbf{全名} \\
\midrule
\endhead

\bottomrule
\endfoot

LLM & Large Language Model (大型語言模型) \\
AI & Artificial Intelligence (人工智慧) \\
EHR & Electronic Health Record (電子病歷) \\
FHIR & Fast Healthcare Interoperability Resources (快速醫療互操作性資源) \\
API & Application Programming Interface (應用程式介面) \\
SR & Success Rate (成功率) \\
MRN & Medical Record Number (病歷號) \\
LOINC & Logical Observation Identifiers Names and Codes \\
SNOMED & Systematized Nomenclature of Medicine \\
NDC & National Drug Code \\
CPT & Current Procedural Terminology \\
PHI & Protected Health Information (受保護健康資訊) \\
CAP & Community-Acquired Pneumonia (社區型肺炎) \\
QA & Question-Answering (問答) \\
\bottomrule
\end{longtable}

\end{document}
